\section{Appendix D -- Version 1.3.3 Provenance and Corpus Boundary}
\label{sec:oc133-provenance-corpus-boundary}

This appendix records what the current manuscript is and what it is not. Core
1.3.3 is a bounded review manuscript built from the public release corpus,
proof sheets, finite checks, comparator material, figures, tables, and
reproducibility companions available for this version. It is not a republication
of an earlier baseline and it is not a claim that every historical source file
has the same evidential status.

\subsection{Purpose}

The main body is written for scientific reading. This appendix is written for
provenance review. It lets a reviewer distinguish between the monograph, the
compact article projection, the methods companion, the reviewer map, and the
machine-readable evidence manifests. Those objects support one another, but
none of them should be mistaken for the whole theory.

\subsection{Current-version boundary}

The present version promotes only bounded claims tied to source and evidence
anchors. Historical material is retained when it explains lineage or provides a
source witness, but it does not automatically promote present-version claims.
When an older statement is broader than the 1.3.3 evidence surface, the public
text narrows it, marks it as illustrative, or moves it into a provenance role.

\subsection{Review questions}

\begin{itemize}
\item Which public corpus object supports the present statement?
\item Which proof sheet, finite check, table, figure, comparator, or replay row
      is being invoked?
\item Is the cited material a promoted support row, an illustrative example, a
      historical source witness, or a planned research route?
\item What finding would reopen the wording?
\end{itemize}

Those questions are the governing boundary of the appendix. They protect the
reader from hidden workflow dependency while keeping internal editorial control
material outside the publication text.
